\documentclass[10pt,a4paper]{amsart}
\usepackage{amsmath,amssymb,amsthm,graphicx}
\usepackage{verbatim}
\newenvironment{reference}{\comment}{\endcomment}
%\newenvironment{reference}{}{}
\newenvironment{slogan}{\comment}{\endcomment}
\newenvironment{history}{\comment}{\endcomment}

% For commutative diagrams we use Xy-pic
\usepackage[all]{xy}

% We use 2cell for 2-commutative diagrams.
\xyoption{2cell}
% Theorem environments.
%
\theoremstyle{plain}
\newtheorem{theorem}[subsection]{Theorem}
\newtheorem{proposition}[subsection]{Proposition}
\newtheorem{lemma}[subsection]{Lemma}

\theoremstyle{definition}
\newtheorem{definition}[subsection]{Definition}
\newtheorem{example}[subsection]{Example}
\newtheorem{exercise}[subsection]{Exercise}
\newtheorem{situation}[subsection]{Situation}

\theoremstyle{remark}
\newtheorem{remark}[subsection]{Remark}
\newtheorem{remarks}[subsection]{Remarks}

\numberwithin{equation}{subsection}

% Macros
%
\def\lim{\mathop{\mathrm{lim}}\nolimits}
\def\colim{\mathop{\mathrm{colim}}\nolimits}
\def\Spec{\mathop{\mathrm{Spec}}}
\def\Hom{\mathop{\mathrm{Hom}}\nolimits}
\def\Ext{\mathop{\mathrm{Ext}}\nolimits}
\def\SheafHom{\mathop{\mathcal{H}\!\mathit{om}}\nolimits}
\def\SheafExt{\mathop{\mathcal{E}\!\mathit{xt}}\nolimits}
\def\Sch{\mathit{Sch}}
\def\Mor{\mathop{\mathrm{Mor}}\nolimits}
\def\Ob{\mathop{\mathrm{Ob}}\nolimits}
\def\Sh{\mathop{\mathit{Sh}}\nolimits}
\def\NL{\mathop{N\!L}\nolimits}
\def\CH{\mathop{\mathrm{CH}}\nolimits}
\def\proetale{{pro\text{-}\acute{e}tale}}
\def\etale{{\acute{e}tale}}
\def\QCoh{\mathit{QCoh}}
\def\Ker{\mathop{\mathrm{Ker}}}
\def\Im{\mathop{\mathrm{Im}}}
\def\Coker{\mathop{\mathrm{Coker}}}
\def\Coim{\mathop{\mathrm{Coim}}}

% Boxtimes
%
\DeclareMathSymbol{\boxtimes}{\mathbin}{AMSa}{"02}

%
% Macros for moduli stacks/spaces
%
\def\QCohstack{\mathcal{QC}\!\mathit{oh}}
\def\Cohstack{\mathcal{C}\!\mathit{oh}}
\def\Spacesstack{\mathcal{S}\!\mathit{paces}}
\def\Quotfunctor{\mathrm{Quot}}
\def\Hilbfunctor{\mathrm{Hilb}}
\def\Curvesstack{\mathcal{C}\!\mathit{urves}}
\def\Polarizedstack{\mathcal{P}\!\mathit{olarized}}
\def\Complexesstack{\mathcal{C}\!\mathit{omplexes}}
% \Pic is the operator that assigns to X its picard group, usage \Pic(X)
% \Picardstack_{X/B} denotes the Picard stack of X over B
% \Picardfunctor_{X/B} denotes the Picard functor of X over B
\def\Pic{\mathop{\mathrm{Pic}}\nolimits}
\def\Picardstack{\mathcal{P}\!\mathit{ic}}
\def\Picardfunctor{\mathrm{Pic}}
\def\Deformationcategory{\mathcal{D}\!\mathit{ef}}
\usepackage{cleveref}
\crefname{originalreference}{original reference}{original references}
\Crefname{originalreference}{Original reference}{Original references}
\makeatletter
\newlabel{lemma-compute-over-affine}{{0CTG}{source}{}{sourcepointer-1}{}}
\newlabel{lemma-compute-over-affine@cref}{{[originalreference][0][]0CTG}{[1][0][]source}{}{}{}}
\makeatother
\crefname{originalreference}{original reference}{original references}
\Crefname{originalreference}{Original reference}{Original references}
\makeatletter
\newlabel{lemma-compute-what-it-should-be}{{0CTD}{source}{}{sourcepointer-2}{}}
\newlabel{lemma-compute-what-it-should-be@cref}{{[originalreference][0][]0CTD}{[1][0][]source}{}{}{}}
\makeatother
\crefname{originalreference}{original reference}{original references}
\Crefname{originalreference}{Original reference}{Original references}
\makeatletter
\newlabel{lemma-relative-pseudo-coherence}{{0CSU}{source}{}{sourcepointer-3}{}}
\newlabel{lemma-relative-pseudo-coherence@cref}{{[originalreference][0][]0CSU}{[1][0][]source}{}{}{}}
\makeatother
\crefname{originalreference}{original reference}{original references}
\Crefname{originalreference}{Original reference}{Original references}
\makeatletter
\newlabel{situation-existence}{{08B7}{source}{}{sourcepointer-4}{}}
\newlabel{situation-existence@cref}{{[originalreference][0][]08B7}{[1][0][]source}{}{}{}}
\makeatother
\title{Near the closed subscheme an open neighbourhood carries a single finite presentation cohomology sheaf}
\author{The Stacks Project Authors}
\date{stacks-project:spaces-flat.tex:local-snapshot-20261012; GNU FDL 1.2 or later}
\begin{document}
\maketitle
\noindent\textit{Editorial note.} All delivered mathematical text is the author's own, copied from the pinned Stacks Project chapter file spaces-flat.tex; no statement or proof step was rewritten, reordered or removed.\par\medskip
\noindent\textit{Editorial note.} The counted proof is the author's complete argument: a reduction to an affine etale neighbourhood, a replacement by a bounded above finite free complex, and an induction trimming the top term by a maximal rank argument\par\medskip
\noindent\textit{Editorial note.} The label lemma-finitely-presented names the counted statement, and the proof environment immediately following it in the source is delivered verbatim as its proof.\par\medskip
\noindent\textit{Editorial note.} The counted statement is carried by the author's own lemma environment in the chapter source.\par\medskip
\noindent\textit{Allowed standard prerequisites (editorial).}\begin{itemize}
\item The situation and the complex K produced in the existence argument, and the canonical map from K to the truncations of the module under consideration
\item Relative pseudo-coherence and the representation of the derived pushforward over an affine etale neighbourhood by a bounded above complex of finite free modules over a smooth presentation
\item The comparison of the cohomology of that complex after base change to the truncation with the corresponding modules, and the behaviour of ranks of maps of finite free modules under principal localisation
\end{itemize}\medskip
\ifdefined\hypertarget\hypertarget{sourcepointer-1}{}\fi
\noindent\textit{Source reference.} Original reference 0CTG: sources/stacks-spaces-flat/pinned-original-source.tex, source lines 2472--2472. Prerequisite pointer invoked at proof line(s) 2544; it names a standard fact of the same chapter that the counted argument uses, and every step of the counted proof built on it is retained verbatim in the delivered body.\par\medskip
\ifdefined\hypertarget\hypertarget{sourcepointer-2}{}\fi
\noindent\textit{Source reference.} Original reference 0CTD: sources/stacks-spaces-flat/pinned-original-source.tex, source lines 2388--2388. Prerequisite pointer invoked at proof line(s) 2510; it names a standard fact of the same chapter that the counted argument uses, and every step of the counted proof built on it is retained verbatim in the delivered body.\par\medskip
\ifdefined\hypertarget\hypertarget{sourcepointer-3}{}\fi
\noindent\textit{Source reference.} Original reference 0CSU: sources/stacks-spaces-flat/pinned-original-source.tex, source lines 2455--2455. Prerequisite pointer invoked at proof line(s) 2537; it names a standard fact of the same chapter that the counted argument uses, and every step of the counted proof built on it is retained verbatim in the delivered body.\par\medskip
\ifdefined\hypertarget\hypertarget{sourcepointer-4}{}\fi
\noindent\textit{Source reference.} Original reference 08B7: sources/stacks-spaces-flat/pinned-original-source.tex, source lines 2362--2362. Prerequisite pointer invoked at proof line(s) 2509; it names a standard fact of the same chapter that the counted argument uses, and every step of the counted proof built on it is retained verbatim in the delivered body.\par\medskip
\begin{lemma}
\label{lemma-finitely-presented}
In Situation \ref{situation-existence} let $K$ be as in
Lemma \ref{lemma-compute-what-it-should-be}.
Denote $X_0 \subset |X|$ the closed subset
consisting of points lying over the closed subset
$\Spec(A_1) = \Spec(A_2) = \ldots$ of $\Spec(A)$.
There exists an open subspace $W \subset X$ containing $X_0$
such that
\begin{enumerate}
\item $H^i(K)|_W$ is zero unless $i = 0$,
\item $\mathcal{F} = H^0(K)|_W$ is of finite presentation, and
\item $\mathcal{F}_n = \mathcal{F} \otimes_{\mathcal{O}_X} \mathcal{O}_{X_n}$.
\end{enumerate}
\end{lemma}

\begin{proof}
Fix $n \geq 1$. By construction there is a canonical map
$K \to \mathcal{F}_n$ in $D_\QCoh(\mathcal{O}_X)$
and hence a canonical map $H^0(K) \to \mathcal{F}_n$
of quasi-coherent sheaves. This explains the meaning of part (3).

\medskip\noindent
Let $x \in X_0$ be a point. We will find an open neighbourhood $W$
of $x$ such that (1), (2), and (3) are true. Since $X_0$ is quasi-compact
this will prove the lemma. Let $U \to X$ be an \'etale morphism
with $U$ affine and $u \in U$ a point mapping to $x$. Since $|U| \to |X|$
is open it suffices to find an open neighbourhood of $u$ in $U$
where (1), (2), and (3) are true. Say $U = \Spec(B)$.
Choose a surjection $P \to B$ with $P$ smooth over $A$.
By Lemma \ref{lemma-relative-pseudo-coherence}
and the definition of relative pseudo-coherence
there exists a bounded above complex $F^\bullet$
of finite free $P$-modules representing
$Ri_*K$ where $i : U \to \Spec(P)$ is the closed
immersion induced by the presentation.
Let $M_n$ be the $B$-module corresponding to $\mathcal{F}_n|_U$.
By Lemma \ref{lemma-compute-over-affine}
$$
H^i(F^\bullet \otimes_A A_n) =
\left\{
\begin{matrix}
0 & \text{if} & i \not = 0 \\
M_n & \text{if} & i = 0
\end{matrix}
\right.
$$
Let $i$ be the maximal index such that $F^i$ is nonzero.
If $i \leq 0$, then (1), (2), and (3) are true.
If not, then $i > 0$ and we see that the rank of the map
$$
F^{i - 1} \to F^i
$$
in the point $u$ is maximal. Hence in an open neighbourhood
of $u$ inside $\Spec(P)$ the rank is maximal. Thus after replacing
$P$ by a principal localization we may assume that the displayed
map is surjective. Since $F^i$ is finite free we may choose
a splitting $F^{i - 1} = F' \oplus F^i$. Then we may
replace $F^\bullet$ by the complex
$$
\ldots \to F^{i - 2} \to F' \to 0 \to \ldots
$$
and we win by induction on $i$.
\end{proof}

\end{document}
